% Relatório da Parte 2 de COS242. Os números e as figuras vêm de tabelas.tex, gerado por
% tabelas.py a partir de data.json (bench/summarize.py).
% Compilar: python3 -I report/parte2/tabelas.py && (cd report/parte2 && tectonic relatorio.tex)
\documentclass[a4paper,11pt]{article}

\usepackage[brazil]{babel}
\usepackage[T1]{fontenc}
\usepackage{libertinus}
\setmonofont{Inconsolatazi4}[Extension=.otf, UprightFont=*-Regular, BoldFont=*-Bold, Scale=0.92]
\usepackage[a4paper,hmargin=2.4cm,top=2.2cm,bottom=2.4cm]{geometry}
\usepackage{microtype}
\usepackage{amsmath}
\usepackage{graphicx}
\usepackage{booktabs}
\usepackage{array}
\usepackage{longtable}
\usepackage[table]{xcolor}
\usepackage{listings}
\usepackage{pgfplots}
\usepackage{caption}
\usepackage{flafter}
\usepackage{placeins}
\usepackage{enumitem}
\usepackage[hidelinks]{hyperref}
\pgfplotsset{compat=1.18}
\usetikzlibrary{arrows.meta,positioning,calc,decorations.pathreplacing}

% ── tipografia e espaçamento ──
\linespread{1.25}
\setlength{\parindent}{1.2em}
\setlength{\parskip}{2pt plus 1pt}
\setlist{nosep,leftmargin=1.5em,topsep=3pt}
\captionsetup{font=small,labelfont=bf,labelsep=period,skip=5pt}
\captionsetup[table]{position=top}
\setlength{\textfloatsep}{12pt plus 2pt minus 3pt}
\setlength{\floatsep}{10pt plus 2pt minus 2pt}
\setlength{\intextsep}{10pt plus 2pt minus 2pt}
\renewcommand{\topfraction}{0.85}
\renewcommand{\bottomfraction}{0.6}
\renewcommand{\textfraction}{0.12}
\renewcommand{\floatpagefraction}{0.75}
\renewcommand{\arraystretch}{1.12}

\definecolor{heapc}{HTML}{1B7F79}
\definecolor{ruim}{HTML}{C0392B}
\definecolor{codkw}{HTML}{1F4E9C}
\definecolor{codtipo}{HTML}{1B7F79}
\definecolor{codcom}{HTML}{7A7A7A}
\lstset{language=C++, basicstyle=\ttfamily\small, columns=fullflexible, keepspaces=true,
  keywordstyle=\color{codkw}\bfseries, commentstyle=\color{codcom},
  morekeywords={uint8,uint32}, keywordstyle=[2]\color{codtipo}\bfseries,
  morekeywords=[2]{Entry}, numbers=none, xleftmargin=1.2em,
  escapeinside={(*}{*)}, frame=none, aboveskip=4pt, belowskip=4pt}
\renewcommand{\lstlistingname}{Listagem}
\newcommand{\cod}[1]{\texttt{#1}}
\newcommand{\W}[1]{W\textsubscript{#1}}
\newcommand{\dist}{\mathit{dist}}
\newcommand{\pai}{\mathit{pai}}
\newcommand{\pos}{\mathit{pos}}
\newcommand{\ini}{\mathit{início}}

\makeatletter
\newcommand{\dado}[1]{\@ifundefined{dado@#1}{\textbf{??#1??}}{\csname dado@#1\endcsname}}
\renewcommand\section{\@startsection{section}{1}{\z@}{-2.2ex plus -.6ex minus -.2ex}%
  {1.0ex plus .2ex}{\normalfont\large\bfseries}}
\renewcommand\subsection{\@startsection{subsection}{2}{\z@}{-1.8ex plus -.5ex minus -.2ex}%
  {0.6ex plus .2ex}{\normalfont\normalsize\bfseries}}
\makeatother
% Gerado por report/parte2/tabelas.py a partir de data.json; não edite à mão.
\pgfplotsset{contadores/.style={yticklabels={{heap, \W{5}},{heap, \W{4}},{heap, \W{3}},{heap, embaralhado},{heap, local},{vetor, \W{5}}}}}
\newcommand{\contadoresEspera}{(77.5,0) (66.4,1) (46.1,2) (66.9,3) (4.2,4) (0.9,5)}
\newcommand{\contadoresResto}{(22.5,0) (33.6,1) (53.9,2) (33.1,3) (95.8,4) (99.1,5)}
\newcommand{\contadoresIPC}{\node[font=\scriptsize, right] at (axis cs:100,0) {0,4};
\node[font=\scriptsize, right] at (axis cs:100,1) {0,7};
\node[font=\scriptsize, right] at (axis cs:100,2) {1,3};
\node[font=\scriptsize, right] at (axis cs:100,3) {0,4};
\node[font=\scriptsize, right] at (axis cs:100,4) {1,3};
\node[font=\scriptsize, right] at (axis cs:100,5) {3,5};}
\newcommand{\contadoresMarcas}{\coordinate (espera0) at (axis cs:77.5,0);
\coordinate (espera1) at (axis cs:66.4,1);
\coordinate (espera2) at (axis cs:46.1,2);
\coordinate (espera3) at (axis cs:66.9,3);
\coordinate (espera4) at (axis cs:4.2,4);
\coordinate (espera5) at (axis cs:0.9,5);}
\newcommand{\contadoresRotulos}{\node[font=\scriptsize, right, inner sep=1.5pt] at (espera0) {77\%};
\node[font=\scriptsize, right, inner sep=1.5pt] at (espera1) {66\%};
\node[font=\scriptsize, right, inner sep=1.5pt] at (espera2) {46\%};
\node[font=\scriptsize, right, inner sep=1.5pt] at (espera3) {67\%};
\node[font=\scriptsize, right, inner sep=1.5pt] at (espera4) {4,2\%};
\node[font=\scriptsize, right, inner sep=1.5pt] at (espera5) {0,9\%};}
\newcommand{\tabelaCaminhos}{\W{3} & 20 & 0,80 & 18 & 10 2858 35293 3905 24445 39345 61279 41321 70540 26049 84450 68904 87862 86778 27694 16608 86355 15760 20\\
 & 30 & 0,89 & 14 & 10 2858 35293 96094 85187 53177 64898 87087 77935 42304 78071 93985 27642 26151 30\\
 & 40 & 0,84 & 10 & 10 2858 35293 3905 24445 86286 14327 19697 15930 46307 40\\
 & 50 & 0,87 & 9 & 10 56663 84437 19040 82510 17622 94992 70948 19878 50\\
 & 60 & 0,97 & 16 & 10 2858 35293 3905 24445 89584 2874 99253 66532 33719 96654 75728 79865 74717 59399 8718 60\\
\midrule
\W{4} & 20 & 2,66 & 13 & 10 532245 877016 554617 430107 437166 805646 841402 663364 655748 634947 318070 807260 20\\
 & 30 & 2,02 & 17 & 10 873980 182382 261305 81769 694525 219466 207765 44636 183747 920911 35338 239369 417194 132461 740528 707581 30\\
 & 40 & 2,52 & 15 & 10 873980 182382 674971 557022 188455 418550 275098 277633 801708 730863 274827 386086 188007 477496 40\\
 & 50 & 2,69 & 18 & 10 873980 182382 261305 81769 694525 382144 748624 663332 149827 349584 735867 9905 716970 472621 407659 994210 630598 50\\
 & 60 & 2,41 & 19 & 10 873980 182382 674971 106317 27985 280367 227336 364950 883474 283156 161217 224210 188021 614364 358029 855615 822462 774634 60\\
\midrule
\W{5} & 20 & 17,34 & 13 & 10 7622850 4615833 6528266 1295458 7103730 8790605 7788327 6766415 7155379 2031914 3908338 8885874 20\\
 & 30 & 17,24 & 18 & 10 2846336 5413988 202621 9270546 6836812 2983314 8775277 5549785 3485759 4341576 8868678 3166991 5117196 7291082 1816279 5664093 8947758 30\\
 & 40 & 17,37 & 14 & 10 1299954 1419794 9885733 5641294 409605 4105715 660832 5850281 3215052 7676193 9008138 4846487 9798152 40\\
 & 50 & 15,43 & 12 & 10 1299954 9927727 8451042 5368621 4276332 3716282 3724752 700154 9271362 6416250 396789 50\\
 & 60 & 17,99 & 22 & 10 1299954 1419794 1899178 7524819 757290 9445351 9611705 9104671 8941517 264365 3945175 3090234 2071841 8424191 2860740 1175531 1925405 7718765 3193827 7824884 1454010 60\\
\midrule
\W{6} & 20 & 11,26 & 20 & 10 7194590 1623776 6004729 5713792 8358296 2533780 5227148 8652974 6001519 3007875 2659882 1975860 7178227 4239765 4333977 5135142 7805172 6937097 2042733 20\\
 & 30 & 10,67 & 22 & 10 6473892 2441538 1548272 787618 5728075 8108429 6007822 5550988 2523255 782814 7306955 153130 9107012 8930108 2199545 9604404 5436447 8094054 8301086 7882916 7776260 30\\
 & 40 & 11,31 & 11 & 10 7194590 1308444 2430454 9006783 9339181 6729573 6733497 5619464 9959710 9649921 40\\
 & 50 & 9,60 & 12 & 10 6473892 2441538 3496708 8719217 3608871 1734087 8196874 3353190 4171489 6448393 8854015 50\\
 & 60 & 10,72 & 20 & 10 6473892 2441538 1548272 787618 9566932 345742 218396 8462300 2279724 9337595 9606830 9928328 8701739 8354591 8135590 4215828 8346899 3411658 6447603 60\\
\midrule
\W{7} & 20 & 93,47 & 33 & 10 44348966 895015 41433396 36342079 15296633 26159915 16943639 38743389 10078159 21570419 11797967 14033688 15395825 37638078 10256175 47812933 16114395 3709139 16167977 42063355 25051663 11357393 1237413 44073100 46623569 16564439 33972685 44128375 38865462 44580311 49167593 34107790 20\\
 & 30 & 95,08 & 32 & 10 44348966 895015 41433396 49215936 31738577 6483504 6934370 37879154 8001897 22183730 7265763 1661582 11612335 1574885 30723220 39141968 27645810 37795681 33920812 34023581 12510615 18298407 31754868 14520632 7334968 721843 7608283 44662249 46198809 161794 37094018 30\\
 & 40 & 100,69 & 46 & 10 44348966 895015 41433396 49215936 31738577 6483504 35576121 13084740 8200352 334647 21020208 33240812 43511183 39829015 31993368 35257514 1449696 2909571 35131499 46634817 8774626 26615224 12496865 25597904 17624904 28515031 9039078 42335617 44883250 25945871 4357404 43686238 2293704 18096524 17747192 14433007 4134285 11543624 49518359 34048991 36125121 20450161 31315434 44798389 37766241 40\\
 & 50 & 95,77 & 34 & 10 44348966 895015 41433396 30012782 5306752 45425007 45928549 3090492 35262022 25302788 39964857 25438408 42769429 18259600 14013095 21407931 29028767 13632589 26269428 19841425 16775149 408668 33679805 18032133 15314852 41061763 21731902 43132195 43482012 248139 24050455 39437117 4560936 50\\
 & 60 & 82,25 & 36 & 10 44348966 895015 41433396 36342079 15296633 26159915 27614435 30636628 46606379 39970228 44691860 28472827 14457897 49527742 23832529 33803765 40354942 41967334 30441838 2384620 9629521 15597152 8376037 19638652 35026345 21270637 16840822 36015359 4125149 20825658 31193177 35236066 37175171 16297527 34154631 60\\}
\newcommand{\tabelaTempos}{\W{3} & 100\,mil & 6,1\,mi & 122,0 & 53,6 $\pm$ 18,3\,ms & 6,55\,s & 122$\times$ & 0,5\\
\W{4} & 1,0\,mi & 31,0\,mi & 62,0 & 888,3 $\pm$ 265,3\,ms & 10,9\,min & 736$\times$ & 1,4\\
\W{5} & 10,0\,mi & 40,0\,mi & 8,0 & 5,48 $\pm$ 1,31\,s & \textit{16,3\,h} & $\sim$11.000$\times$ & 4,7\\
\W{6} & 10,0\,mi & 73,4\,mi & 14,7 & 7,19 $\pm$ 0,40\,s & \textit{16,2\,h} & $\sim$8.100$\times$ & 3,7\\
\W{7} & 50,0\,mi & 60,0\,mi & 2,4 & 31,19 $\pm$ 7,65\,s & \textit{16,0\,dias} & $\sim$44.000$\times$ & 11,1\\}
\newcommand{\latencia}{(16384,1.109) (24576,1.115) (32768,1.117) (49152,1.105) (65536,1.114) (98304,1.104) (131072,1.346) (196608,4.166) (262144,4.929) (393216,5.761) (524288,6.096) (786432,6.507) (1048576,6.512) (1572864,6.935) (2097152,6.801) (3145728,7.336) (4194304,8.504) (6291456,10.315) (8388608,12.54) (12582912,50.962) (16777216,49.437) (25165824,75.719) (33554432,87.815) (50331648,99.354) (67108864,104.201) (100663296,108.778) (134217728,114.129) (201326592,115.031) (268435456,115.753) (402653184,115.727) (536870912,118.824) (805306368,121.134) (1073741824,121.188) (1610612736,121.897) (2147483648,137.135)}
\newcommand{\marcasW}{\draw (axis cs:800000,0.7) -- (axis cs:800000,0.85) node[above, font=\tiny, inner sep=1pt] {\W{3}};
\draw (axis cs:8000000,0.7) -- (axis cs:8000000,0.85) node[above, font=\tiny, inner sep=1pt] {\W{4}};
\draw (axis cs:80000000,0.7) -- (axis cs:80000000,0.85) node[above, font=\tiny, inner sep=1pt] {\W{5}, \W{6}};
\draw (axis cs:400000000,0.7) -- (axis cs:400000000,0.85) node[above, font=\tiny, inner sep=1pt] {\W{7}};}
\pgfplotsset{grafosW/.style={symbolic x coords={W_3,W_4,W_5,W_6,W_7}, xticklabels={\W{3},\W{4},\W{5},\W{6},\W{7}}}}
\newcommand{\barrasHeap}{(W_3,0.0535502) (W_4,0.888338) (W_5,5.47735) (W_6,7.1858) (W_7,31.1917)}
\newcommand{\barrasVetor}{(W_3,6.54866) (W_4,653.446)}
\newcommand{\barrasVetorProj}{(W_5,58744) (W_6,58398) (W_7,1.38635e+06)}
\newcommand{\razoesW}{\node[font=\scriptsize, above, inner sep=2pt] at (axis cs:W_3,6.54866) [xshift=5pt] {122$\times$};
\node[font=\scriptsize, above, inner sep=2pt] at (axis cs:W_4,653.446) [xshift=5pt] {736$\times$};
\node[font=\scriptsize, above, inner sep=2pt] at (axis cs:W_5,58744) [xshift=5pt] {$\sim$11.000$\times$};
\node[font=\scriptsize, above, inner sep=2pt] at (axis cs:W_6,58398) [xshift=5pt] {$\sim$8.100$\times$};
\node[font=\scriptsize, above, inner sep=2pt] at (axis cs:W_7,1.38635e+06) [xshift=5pt] {$\sim$44.000$\times$};}
\newcommand{\pontosEscala}{(102978790.97272025,0.0535502) (637809755.723865,0.888338) (1162674716.9430935,5.47735) (1939420032.5859938,7.1858) (2813296697.925454,31.1917)}
\newcommand{\retaEscala}{(51489395.486360125,0.0267751) (5626593395.850908,2.92589)}
\newcommand{\rotulosEscala}{\node[font=\scriptsize, anchor=north west, inner sep=2pt] at (axis cs:1.02979e+08,0.0535502) {\W{3}};
\node[font=\scriptsize, anchor=north west, inner sep=2pt] at (axis cs:6.3781e+08,0.888338) {\W{4}};
\node[font=\scriptsize, anchor=north west, inner sep=2pt] at (axis cs:1.16267e+09,5.47735) {\W{5}};
\node[font=\scriptsize, anchor=north west, inner sep=2pt] at (axis cs:1.93942e+09,7.1858) {\W{6}};
\node[font=\scriptsize, anchor=north west, inner sep=2pt] at (axis cs:2.8133e+09,31.1917) {\W{7}};}
\newcommand{\pontosVetorPos}{(100000,0.58068) (1000000,0.58948) (10000000,0.58744) (10000000,0.58398) (50000000,0.55454)}
\newcommand{\pontosAtualizacoes}{(121.99882,4.37223) (61.999956,3.71408) (7.999999,1.83981) (14.6806744,2.33595) (2.39999996,1.07483)}
\newcommand{\rotulosAtualizacoes}{\node[font=\scriptsize, anchor=north west, inner sep=2pt] at (axis cs:121.999,4.37223) {\W{3}};
\node[font=\scriptsize, anchor=north west, inner sep=2pt] at (axis cs:62,3.71408) {\W{4}};
\node[font=\scriptsize, anchor=north west, inner sep=2pt] at (axis cs:8,1.83981) {\W{5}};
\node[font=\scriptsize, anchor=north west, inner sep=2pt] at (axis cs:14.6807,2.33595) {\W{6}};
\node[font=\scriptsize, anchor=north west, inner sep=2pt] at (axis cs:2.4,1.07483) {\W{7}};}
\expandafter\def\csname heapTam@W_3\endcsname{(1e-05,0) (0.01001,0.69481) (0.02001,0.89377) (0.03001,0.94442) (0.04001,0.95252) (0.05001,0.94774) (0.06001,0.93929) (0.07001,0.92973) (0.08001,0.91986) (0.09001,0.90996) (0.10001,0.89999) (0.11001,0.88999) (0.12001,0.87999) (0.13001,0.86999) (0.14001,0.85999) (0.15001,0.84999) (0.16001,0.83999) (0.17001,0.82999) (0.18001,0.81999) (0.19001,0.80999) (0.20001,0.79999) (0.21001,0.78999) (0.22001,0.77999) (0.23001,0.76999) (0.24001,0.75999) (0.25001,0.74999) (0.26001,0.73999) (0.27001,0.72999) (0.28001,0.71999) (0.29001,0.70999) (0.30001,0.69999) (0.31001,0.68999) (0.32001,0.67999) (0.33001,0.66999) (0.34001,0.65999) (0.35001,0.64999) (0.36001,0.63999) (0.37001,0.62999) (0.38001,0.61999) (0.39001,0.60999) (0.40001,0.59999) (0.41001,0.58999) (0.42001,0.57999) (0.43001,0.56999) (0.44001,0.55999) (0.45001,0.54999) (0.46001,0.53999) (0.47001,0.52999) (0.48001,0.51999) (0.49001,0.50999) (0.50001,0.49999) (0.51001,0.48999) (0.52001,0.47999) (0.53001,0.46999) (0.54001,0.45999) (0.55001,0.44999) (0.56001,0.43999) (0.57001,0.42999) (0.58001,0.41999) (0.59001,0.40999) (0.60001,0.39999) (0.61001,0.38999) (0.62001,0.37999) (0.63001,0.36999) (0.64001,0.35999) (0.65001,0.34999) (0.66001,0.33999) (0.67001,0.32999) (0.68001,0.31999) (0.69001,0.30999) (0.70001,0.29999) (0.71001,0.28999) (0.72001,0.27999) (0.73001,0.26999) (0.74001,0.25999) (0.75001,0.24999) (0.76001,0.23999) (0.77001,0.22999) (0.78001,0.21999) (0.79001,0.20999) (0.80001,0.19999) (0.81001,0.18999) (0.82001,0.17999) (0.83001,0.16999) (0.84001,0.15999) (0.85001,0.14999) (0.86001,0.13999) (0.87001,0.12999) (0.88001,0.11999) (0.89001,0.10999) (0.90001,0.09999) (0.91001,0.08999) (0.92001,0.07999) (0.93001,0.06999) (0.94001,0.05999) (0.95001,0.04999) (0.96001,0.03999) (0.97001,0.02999) (0.98001,0.01999) (0.99001,0.00999) (0.99801,0.00199)}
\expandafter\def\csname heapTam@W_4\endcsname{(1e-06,0) (0.010001,0.451817) (0.020001,0.690943) (0.030001,0.814723) (0.040001,0.876355) (0.050001,0.905086) (0.060001,0.915918) (0.070001,0.917145) (0.080001,0.913084) (0.090001,0.906248) (0.100001,0.898047) (0.110001,0.888964) (0.120001,0.879414) (0.130001,0.869687) (0.140001,0.859822) (0.150001,0.849903) (0.160001,0.839948) (0.170001,0.829967) (0.180001,0.819982) (0.190001,0.809994) (0.200001,0.799997) (0.210001,0.789997) (0.220001,0.779999) (0.230001,0.769999) (0.240001,0.759999) (0.250001,0.749999) (0.260001,0.739999) (0.270001,0.729999) (0.280001,0.719999) (0.290001,0.709999) (0.300001,0.699999) (0.310001,0.689999) (0.320001,0.679999) (0.330001,0.669999) (0.340001,0.659999) (0.350001,0.649999) (0.360001,0.639999) (0.370001,0.629999) (0.380001,0.619999) (0.390001,0.609999) (0.400001,0.599999) (0.410001,0.589999) (0.420001,0.579999) (0.430001,0.569999) (0.440001,0.559999) (0.450001,0.549999) (0.460001,0.539999) (0.470001,0.529999) (0.480001,0.519999) (0.490001,0.509999) (0.500001,0.499999) (0.510001,0.489999) (0.520001,0.479999) (0.530001,0.469999) (0.540001,0.459999) (0.550001,0.449999) (0.560001,0.439999) (0.570001,0.429999) (0.580001,0.419999) (0.590001,0.409999) (0.600001,0.399999) (0.610001,0.389999) (0.620001,0.379999) (0.630001,0.369999) (0.640001,0.359999) (0.650001,0.349999) (0.660001,0.339999) (0.670001,0.329999) (0.680001,0.319999) (0.690001,0.309999) (0.700001,0.299999) (0.710001,0.289999) (0.720001,0.279999) (0.730001,0.269999) (0.740001,0.259999) (0.750001,0.249999) (0.760001,0.239999) (0.770001,0.229999) (0.780001,0.219999) (0.790001,0.209999) (0.800001,0.199999) (0.810001,0.189999) (0.820001,0.179999) (0.830001,0.169999) (0.840001,0.159999) (0.850001,0.149999) (0.860001,0.139999) (0.870001,0.129999) (0.880001,0.119999) (0.890001,0.109999) (0.900001,0.099999) (0.910001,0.089999) (0.920001,0.079999) (0.930001,0.069999) (0.940001,0.059999) (0.950001,0.049999) (0.960001,0.039999) (0.970001,0.029999) (0.980001,0.019999) (0.990001,0.009999) (0.998001,0.001999)}
\expandafter\def\csname heapTam@W_5\endcsname{(1e-07,0) (0.0100001,0.06377) (0.0200001,0.122124) (0.0300001,0.175612) (0.0400001,0.224508) (0.0500001,0.269163) (0.0600001,0.309786) (0.0700001,0.346737) (0.0800001,0.380318) (0.0900001,0.410812) (0.1000001,0.43828) (0.1100001,0.463069) (0.1200001,0.485303) (0.1300001,0.505135) (0.1400001,0.522836) (0.1500001,0.538503) (0.1600001,0.552299) (0.1700001,0.564278) (0.1800001,0.574667) (0.1900001,0.583585) (0.2000001,0.591073) (0.2100001,0.597218) (0.2200001,0.602236) (0.2300001,0.606055) (0.2400001,0.608912) (0.2500001,0.610741) (0.2600001,0.611645) (0.2700001,0.611822) (0.2800001,0.611138) (0.2900001,0.609764) (0.3000001,0.607797) (0.3100001,0.60516) (0.3200001,0.601993) (0.3300001,0.598266) (0.3400001,0.594084) (0.3500001,0.589451) (0.3600001,0.584437) (0.3700001,0.578987) (0.3800001,0.573159) (0.3900001,0.567037) (0.4000001,0.560634) (0.4100001,0.553901) (0.4200001,0.546961) (0.4300001,0.539728) (0.4400001,0.532309) (0.4500001,0.524669) (0.4600001,0.516849) (0.4700001,0.50885) (0.4800001,0.500694) (0.4900001,0.492365) (0.5000001,0.483893) (0.5100001,0.475304) (0.5200001,0.46661) (0.5300001,0.457799) (0.5400001,0.448915) (0.5500001,0.439913) (0.5600001,0.430836) (0.5700001,0.421682) (0.5800001,0.412473) (0.5900001,0.403169) (0.6000001,0.393813) (0.6100001,0.384416) (0.6200001,0.374958) (0.6300001,0.365446) (0.6400001,0.355905) (0.6500001,0.346311) (0.6600001,0.33668) (0.6700001,0.327017) (0.6800001,0.317329) (0.6900001,0.30761) (0.7000001,0.297865) (0.7100001,0.2881) (0.7200001,0.2783) (0.7300001,0.268485) (0.7400001,0.258662) (0.7500001,0.248819) (0.7600001,0.238954) (0.7700001,0.229076) (0.7800001,0.219185) (0.7900001,0.209283) (0.8000001,0.199371) (0.8100001,0.189449) (0.8200001,0.179521) (0.8300001,0.169588) (0.8400001,0.159648) (0.8500001,0.149698) (0.8600001,0.139741) (0.8700001,0.129782) (0.8800001,0.119817) (0.8900001,0.109849) (0.9000001,0.0998736) (0.9100001,0.0898962) (0.9200001,0.0799175) (0.9300001,0.0699362) (0.9400001,0.059953) (0.9500001,0.049966) (0.9600001,0.0399762) (0.9700001,0.0299855) (0.9800001,0.0199923) (0.9900001,0.0099974) (0.9980001,0.0019997)}
\expandafter\def\csname heapTam@W_6\endcsname{(1e-07,0) (0.0099996,0.126414) (0.0199991,0.234309) (0.0299986,0.326354) (0.0399981,0.404102) (0.0499976,0.469964) (0.0599971,0.525361) (0.0699966,0.572058) (0.0799961,0.610837) (0.0899956,0.643045) (0.0999951,0.669398) (0.1099946,0.69094) (0.1199941,0.708174) (0.1299936,0.721609) (0.1399931,0.731886) (0.1499926,0.739352) (0.1599921,0.744467) (0.1699916,0.74746) (0.1799911,0.748749) (0.1899906,0.748463) (0.1999901,0.746868) (0.2099896,0.744116) (0.2199891,0.740393) (0.2299886,0.735825) (0.2399881,0.730519) (0.2499876,0.724515) (0.2599871,0.717986) (0.2699866,0.710982) (0.2799861,0.703593) (0.2899856,0.69585) (0.2999851,0.687762) (0.3099846,0.679436) (0.3199841,0.670867) (0.3299836,0.662106) (0.3399831,0.653186) (0.3499826,0.644115) (0.3599821,0.634913) (0.3699816,0.625606) (0.3799811,0.616207) (0.3899806,0.606733) (0.3999801,0.59718) (0.4099796,0.587566) (0.4199791,0.577895) (0.4299786,0.568191) (0.4399781,0.558438) (0.4499776,0.548652) (0.4599771,0.538831) (0.4699766,0.528992) (0.4799761,0.519136) (0.4899756,0.50925) (0.4999751,0.499357) (0.5099746,0.489448) (0.5199741,0.479524) (0.5299736,0.469594) (0.5399731,0.459653) (0.5499726,0.449707) (0.5599721,0.439749) (0.5699716,0.429785) (0.5799711,0.419816) (0.5899706,0.409843) (0.5999701,0.399868) (0.6099696,0.389892) (0.6199691,0.379912) (0.6299686,0.369929) (0.6399681,0.359944) (0.6499676,0.349957) (0.6599671,0.339966) (0.6699666,0.329975) (0.6799661,0.319984) (0.6899656,0.309992) (0.6999651,0.299998) (0.7099646,0.290003) (0.7199641,0.280008) (0.7299636,0.270012) (0.7399631,0.260016) (0.7499626,0.250018) (0.7599621,0.240021) (0.7699616,0.230023) (0.7799611,0.220026) (0.7899606,0.210028) (0.7999601,0.200029) (0.8099596,0.190031) (0.8199591,0.180033) (0.8299586,0.170035) (0.8399581,0.160036) (0.8499576,0.150037) (0.8599571,0.140039) (0.8699566,0.13004) (0.8799561,0.120041) (0.8899556,0.110042) (0.8999551,0.100043) (0.9099546,0.0900435) (0.9199541,0.0800443) (0.9299536,0.070045) (0.9399531,0.0600457) (0.9499526,0.0500463) (0.9599521,0.0400472) (0.9699516,0.0300479) (0.9799511,0.0200484) (0.9899506,0.010049) (0.9999501,4.96e-05)}
\expandafter\def\csname heapTam@W_7\endcsname{(2e-08,0) (0.01000002,0.00585598) (0.02000002,0.0115759) (0.03000002,0.0171522) (0.04000002,0.0225846) (0.05000002,0.0278843) (0.06000002,0.0330363) (0.07000002,0.0380474) (0.08000002,0.0429391) (0.09000002,0.0476966) (0.10000002,0.0523114) (0.11000002,0.0568061) (0.12000002,0.0611738) (0.13000002,0.0654037) (0.14000002,0.0695033) (0.15000002,0.0734617) (0.16000002,0.077289) (0.17000002,0.0810176) (0.18000002,0.0845888) (0.19000002,0.088048) (0.20000002,0.0913775) (0.21000002,0.0945787) (0.22000002,0.0976607) (0.23000002,0.10063) (0.24000002,0.103479) (0.25000002,0.106202) (0.26000002,0.10878) (0.27000002,0.111243) (0.28000002,0.113581) (0.29000002,0.11579) (0.30000002,0.117893) (0.31000002,0.119876) (0.32000002,0.121753) (0.33000002,0.123509) (0.34000002,0.125133) (0.35000002,0.126651) (0.36000002,0.128049) (0.37000002,0.129349) (0.38000002,0.13053) (0.39000002,0.131608) (0.40000002,0.132574) (0.41000002,0.133405) (0.42000002,0.134109) (0.43000002,0.13471) (0.44000002,0.135237) (0.45000002,0.135652) (0.46000002,0.135944) (0.47000002,0.136119) (0.48000002,0.136201) (0.49000002,0.136156) (0.50000002,0.136018) (0.51000002,0.135774) (0.52000002,0.135437) (0.53000002,0.135004) (0.54000002,0.13444) (0.55000002,0.133784) (0.56000002,0.133044) (0.57000002,0.132186) (0.58000002,0.131224) (0.59000002,0.130168) (0.60000002,0.129005) (0.61000002,0.127721) (0.62000002,0.126349) (0.63000002,0.124888) (0.64000002,0.123308) (0.65000002,0.121617) (0.66000002,0.119843) (0.67000002,0.117968) (0.68000002,0.116001) (0.69000002,0.113914) (0.70000002,0.111746) (0.71000002,0.109499) (0.72000002,0.107114) (0.73000002,0.104632) (0.74000002,0.10208) (0.75000002,0.099415) (0.76000002,0.096649) (0.77000002,0.093795) (0.78000002,0.0908606) (0.79000002,0.0878032) (0.80000002,0.0846696) (0.81000002,0.0814164) (0.82000002,0.0780781) (0.83000002,0.0746326) (0.84000002,0.0710893) (0.85000002,0.0674542) (0.86000002,0.063719) (0.87000002,0.0598868) (0.88000002,0.0559458) (0.89000002,0.0518974) (0.90000002,0.0477343) (0.91000002,0.0434833) (0.92000002,0.0391398) (0.93000002,0.0346714) (0.94000002,0.0300867) (0.95000002,0.0254017) (0.96000002,0.0205933) (0.97000002,0.0156657) (0.98000002,0.0105971) (0.99000002,0.00539128) (0.99800002,0.00109756)}
\newcommand{\arvoreColab}{\node[raiz] (n31) at (6.30,0.00) {Edsger W. Dijkstra};
\node[pessoa] (n18) at (2.10,-0.92) {A. J. M. van Gasteren};
\node[pessoa] (n30) at (10.50,-0.92) {John R. Rice};
\node[pessoa] (n17) at (2.10,-1.84) {Gerard Tel};
\node[pessoa] (n29) at (10.50,-1.84) {Dan C. Marinescu};
\node[pessoa] (n16) at (2.10,-2.76) {Hans L. Bodlaender};
\node[pessoa] (n23) at (8.40,-2.76) {Chuang Lin};
\node[pessoa] (n28) at (12.60,-2.76) {Howard Jay Siegel};
\node[pessoa] (n8) at (0.00,-3.68) {Dimitrios M. Thilikos};
\node[pessoa] (n15) at (4.20,-3.68) {Jan van Leeuwen};
\node[pessoa] (n22) at (8.40,-3.68) {Bo Li};
\node[pessoa] (n27) at (12.60,-3.68) {Edwin K. P. Chong};
\node[pessoa] (n7) at (0.00,-4.60) {Prabhakar Ragde};
\node[pessoa] (n14) at (4.20,-4.60) {Mark H. Overmars};
\node[pessoa] (n21) at (8.40,-4.60) {Y. Thomas Hou};
\node[pessoa] (n26) at (12.60,-4.60) {Ness B. Shroff};
\node[pessoa] (n6) at (0.00,-5.52) {Avi Wigderson};
\node[pessoa] (n13) at (4.20,-5.52) {Micha Sharir};
\node[pessoa] (n20) at (8.40,-5.52) {Zhi-Li Zhang};
\node[pessoa] (n25) at (12.60,-5.52) {R. Srikant};
\node[pessoa] (n5) at (0.00,-6.44) {Eli Upfal};
\node[pessoa] (n12) at (4.20,-6.44) {Haim Kaplan};
\node[pessoa] (n19) at (8.40,-6.44) {Donald F. Towsley};
\node[pessoa] (n24) at (12.60,-6.44) {Albert G. Greenberg};
\node[folha] (n2) at (8.40,-7.58) {Daniel R. Figueiredo\\\textbf{2,94}};
\node[folha] (n3) at (12.60,-7.58) {J. B. Kruskal\\\textbf{3,48}};
\node[pessoa] (n4) at (0.00,-7.36) {Prabhakar Raghavan};
\node[pessoa] (n11) at (4.20,-7.36) {Robert Endre Tarjan};
\node[folha] (n0) at (0.00,-8.50) {Jon M. Kleinberg\\\textbf{2,71}};
\node[pessoa] (n10) at (4.20,-8.28) {Andrew V. Goldberg};
\node[pessoa] (n9) at (4.20,-9.20) {Serge A. Plotkin};
\node[folha] (n1) at (4.20,-10.34) {Éva Tardos\\\textbf{2,75}};
\draw[barato] (n31) -- node[peso] {0,40} (n18);
\draw[caro] (n31) -- node[peso] {1,00} (n30);
\draw[caro] (n18) -- node[peso] {1,00} (n17);
\draw[barato] (n17) -- node[peso] {0,31} (n16);
\draw[barato] (n16) -- node[peso] {0,10} (n8);
\draw[barato] (n16) -- node[peso] {0,17} (n15);
\draw[barato] (n8) -- node[peso] {0,21} (n7);
\draw[barato] (n7) -- node[peso] {0,27} (n6);
\draw[barato] (n6) -- node[peso] {0,18} (n5);
\draw[barato] (n5) -- node[peso] {0,10} (n4);
\draw[barato] (n4) -- node[peso] {0,14} (n0);
\draw[barato] (n15) -- node[peso] {0,11} (n14);
\draw[barato] (n14) -- node[peso] {0,11} (n13);
\draw[barato] (n13) -- node[peso] {0,14} (n12);
\draw[barato] (n12) -- node[peso] {0,10} (n11);
\draw[barato] (n11) -- node[peso] {0,13} (n10);
\draw[barato] (n10) -- node[peso] {0,17} (n9);
\draw[barato] (n9) -- node[peso] {0,13} (n1);
\draw[barato] (n30) -- node[peso] {0,13} (n29);
\draw[barato] (n29) -- node[peso] {0,30} (n23);
\draw[barato] (n29) -- node[peso] {0,30} (n28);
\draw[barato] (n23) -- node[peso] {0,39} (n22);
\draw[barato] (n22) -- node[peso] {0,26} (n21);
\draw[barato] (n21) -- node[peso] {0,22} (n20);
\draw[barato] (n20) -- node[peso] {0,17} (n19);
\draw[barato] (n19) -- node[peso] {0,47} (n2);
\draw[barato] (n28) -- node[peso] {0,26} (n27);
\draw[barato] (n27) -- node[peso] {0,10} (n26);
\draw[barato] (n26) -- node[peso] {0,29} (n25);
\draw[barato] (n25) -- node[peso] {0,40} (n24);
\draw[caro] (n24) -- node[peso] {1,00} (n3);}
\newcommand{\tabelaDistancias}{20 & 0,80 & 18 & 2,66 & 13 & 17,34 & 13 & 11,26 & 20 & 93,47 & 33\\
30 & 0,89 & 14 & 2,02 & 17 & 17,24 & 18 & 10,67 & 22 & 95,08 & 32\\
40 & 0,84 & 10 & 2,52 & 15 & 17,37 & 14 & 11,31 & 11 & 100,69 & 46\\
50 & 0,87 & 9 & 2,69 & 18 & 15,43 & 12 & 9,60 & 12 & 95,77 & 34\\
60 & 0,97 & 16 & 2,41 & 19 & 17,99 & 22 & 10,72 & 20 & 82,25 & 36\\}
\expandafter\def\csname dado@k\endcsname{100}
\expandafter\def\csname dado@seed\endcsname{2026}
\expandafter\def\csname dado@cpu\endcsname{Apple M3 Pro}
\expandafter\def\csname dado@ram\endcsname{36\,GiB}
\expandafter\def\csname dado@compiler\endcsname{Apple clang version 21.0.0}
\expandafter\def\csname dado@ghz\endcsname{3,58}
\expandafter\def\csname dado@scan.lib\endcsname{0,58}
\expandafter\def\csname dado@scan.lib.min\endcsname{0,55}
\expandafter\def\csname dado@scan.lib.max\endcsname{0,63}
\expandafter\def\csname dado@scan.lib.cycles\endcsname{2,1}
\expandafter\def\csname dado@scan.chains\endcsname{0,26}
\expandafter\def\csname dado@scan.chains.x\endcsname{2,3$\times$}
\expandafter\def\csname dado@scan.simd\endcsname{0,09}
\expandafter\def\csname dado@scan.simd.x\endcsname{6,8$\times$}
\expandafter\def\csname dado@scan.simd.big\endcsname{0,14}
\expandafter\def\csname dado@scan.big.n\endcsname{50,0\,mi}
\expandafter\def\csname dado@scan.simd.gbs\endcsname{56}
\expandafter\def\csname dado@pops.min\endcsname{0,55}
\expandafter\def\csname dado@pops.max\endcsname{0,59}
\expandafter\def\csname dado@pops.cycles\endcsname{2,1}
\expandafter\def\csname dado@val.W_3.err\endcsname{1,1\%}
\expandafter\def\csname dado@val.W_3.runs\endcsname{100}
\expandafter\def\csname dado@val.W_3.slowest\endcsname{3,5$\times$}
\expandafter\def\csname dado@val.W_3.spread\endcsname{268\%}
\expandafter\def\csname dado@val.W_4.err\endcsname{$-$10,6\%}
\expandafter\def\csname dado@val.W_4.runs\endcsname{5}
\expandafter\def\csname dado@val.W_4.slowest\endcsname{1,1$\times$}
\expandafter\def\csname dado@val.W_4.spread\endcsname{23\%}
\expandafter\def\csname dado@pops.limit.W_7\endcsname{300}
\expandafter\def\csname dado@pops.limit.W_5\endcsname{1.500}
\expandafter\def\csname dado@lat.l1\endcsname{1,1}
\expandafter\def\csname dado@lat.l2\endcsname{8,5}
\expandafter\def\csname dado@lat.dram\endcsname{121}
\expandafter\def\csname dado@layout.5.local\endcsname{0,82}
\expandafter\def\csname dado@layout.5.random\endcsname{0,97}
\expandafter\def\csname dado@layout.5.x\endcsname{1,2$\times$}
\expandafter\def\csname dado@layout.6.local\endcsname{0,71}
\expandafter\def\csname dado@layout.6.random\endcsname{1,72}
\expandafter\def\csname dado@layout.6.x\endcsname{2,4$\times$}
\expandafter\def\csname dado@layout.7.local\endcsname{0,60}
\expandafter\def\csname dado@layout.7.random\endcsname{1,85}
\expandafter\def\csname dado@layout.7.x\endcsname{3,1$\times$}
\expandafter\def\csname dado@layout.max.local\endcsname{0,55}
\expandafter\def\csname dado@layout.max.random\endcsname{1,79}
\expandafter\def\csname dado@layout.max.x\endcsname{3,3$\times$}
\expandafter\def\csname dado@layout.maxn\endcsname{50,0\,mi}
\expandafter\def\csname dado@order.W_3.gap\endcsname{0,333}
\expandafter\def\csname dado@order.W_3.random\endcsname{0,333}
\expandafter\def\csname dado@order.W_3.bfs\endcsname{0,329}
\expandafter\def\csname dado@order.W_3.bfs.change\endcsname{$-$15\%}
\expandafter\def\csname dado@order.W_3.random.change\endcsname{$-$2\%}
\expandafter\def\csname dado@order.W_3.pushes\endcsname{4,4}
\expandafter\def\csname dado@order.W_3.peak\endcsname{95\,mil}
\expandafter\def\csname dado@order.W_3.peak.pct\endcsname{95\%}
\expandafter\def\csname dado@order.W_5.gap\endcsname{0,333}
\expandafter\def\csname dado@order.W_5.random\endcsname{0,333}
\expandafter\def\csname dado@order.W_5.bfs\endcsname{0,272}
\expandafter\def\csname dado@order.W_5.bfs.change\endcsname{$-$3\%}
\expandafter\def\csname dado@order.W_5.random.change\endcsname{$-$3\%}
\expandafter\def\csname dado@order.W_5.pushes\endcsname{1,8}
\expandafter\def\csname dado@order.W_5.peak\endcsname{6,1\,mi}
\expandafter\def\csname dado@order.W_5.peak.pct\endcsname{61\%}
\expandafter\def\csname dado@layout.peak.max\endcsname{425}
\expandafter\def\csname dado@layout.peak.kb\endcsname{7\,kB}
\expandafter\def\csname dado@ctr.layout.n\endcsname{10,0\,mi}
\expandafter\def\csname dado@ctr.layout.instr\endcsname{1\%}
\expandafter\def\csname dado@ctr.layout.cycles.x\endcsname{2,9$\times$}
\expandafter\def\csname dado@ctr.layout.l1.x\endcsname{12$\times$}
\expandafter\def\csname dado@ctr.layout.tlb.x\endcsname{48$\times$}
\expandafter\def\csname dado@ctr.local.wait\endcsname{4,2\%}
\expandafter\def\csname dado@ctr.random.wait\endcsname{67\%}
\expandafter\def\csname dado@ctr.vec.wait\endcsname{0,9\%}
\expandafter\def\csname dado@ctr.vec.ipc\endcsname{3,5}
\expandafter\def\csname dado@ctr.vec.cpe\endcsname{2,0}
\expandafter\def\csname dado@ctr.W_3.wait\endcsname{46\%}
\expandafter\def\csname dado@ctr.W_3.ipc\endcsname{1,3}
\expandafter\def\csname dado@ctr.W_3.tlb\endcsname{2}
\expandafter\def\csname dado@ctr.W_4.wait\endcsname{66\%}
\expandafter\def\csname dado@ctr.W_4.ipc\endcsname{0,7}
\expandafter\def\csname dado@ctr.W_4.tlb\endcsname{4}
\expandafter\def\csname dado@ctr.W_5.wait\endcsname{77\%}
\expandafter\def\csname dado@ctr.W_5.ipc\endcsname{0,4}
\expandafter\def\csname dado@ctr.W_5.tlb\endcsname{29}
\expandafter\def\csname dado@W_3.heap\endcsname{53,6\,ms}
\expandafter\def\csname dado@W_3.n\endcsname{100\,mil}
\expandafter\def\csname dado@W_3.heapk\endcsname{5,36\,s}
\expandafter\def\csname dado@W_3.vec\endcsname{6,55\,s}
\expandafter\def\csname dado@W_3.vecest\endcsname{5,81\,s}
\expandafter\def\csname dado@W_3.veck\endcsname{9,7\,min}
\expandafter\def\csname dado@W_3.ratio\endcsname{122$\times$}
\expandafter\def\csname dado@W_3.heapns\endcsname{0,5}
\expandafter\def\csname dado@W_3.degree\endcsname{122,0}
\expandafter\def\csname dado@W_3.csr\endcsname{147\,MB}
\expandafter\def\csname dado@W_3.pico\endcsname{249\,MB}
\expandafter\def\csname dado@W_3.matriz\endcsname{90,0\,GB}
\expandafter\def\csname dado@W_3.chainsest\endcsname{2,59\,s}
\expandafter\def\csname dado@W_3.chainsratio\endcsname{48$\times$}
\expandafter\def\csname dado@W_4.heap\endcsname{888,3\,ms}
\expandafter\def\csname dado@W_4.n\endcsname{1,0\,mi}
\expandafter\def\csname dado@W_4.heapk\endcsname{88,83\,s}
\expandafter\def\csname dado@W_4.vec\endcsname{10,9\,min}
\expandafter\def\csname dado@W_4.vecest\endcsname{9,8\,min}
\expandafter\def\csname dado@W_4.veck\endcsname{16,4\,h}
\expandafter\def\csname dado@W_4.ratio\endcsname{736$\times$}
\expandafter\def\csname dado@W_4.heapns\endcsname{1,4}
\expandafter\def\csname dado@W_4.degree\endcsname{62,0}
\expandafter\def\csname dado@W_4.csr\endcsname{748\,MB}
\expandafter\def\csname dado@W_4.pico\endcsname{1,3\,GB}
\expandafter\def\csname dado@W_4.matriz\endcsname{9,0\,TB}
\expandafter\def\csname dado@W_4.chainsest\endcsname{4,4\,min}
\expandafter\def\csname dado@W_4.chainsratio\endcsname{296$\times$}
\expandafter\def\csname dado@W_5.heap\endcsname{5,48\,s}
\expandafter\def\csname dado@W_5.n\endcsname{10,0\,mi}
\expandafter\def\csname dado@W_5.heapk\endcsname{9,1\,min}
\expandafter\def\csname dado@W_5.vecest\endcsname{16,3\,h}
\expandafter\def\csname dado@W_5.veck\endcsname{68,0\,dias}
\expandafter\def\csname dado@W_5.ratio\endcsname{$\sim$11.000$\times$}
\expandafter\def\csname dado@W_5.heapns\endcsname{4,7}
\expandafter\def\csname dado@W_5.degree\endcsname{8,0}
\expandafter\def\csname dado@W_5.csr\endcsname{1,0\,GB}
\expandafter\def\csname dado@W_5.pico\endcsname{1,9\,GB}
\expandafter\def\csname dado@W_5.matriz\endcsname{900,0\,TB}
\expandafter\def\csname dado@W_5.chainsest\endcsname{7,4\,h}
\expandafter\def\csname dado@W_5.chainsratio\endcsname{$\sim$4.900$\times$}
\expandafter\def\csname dado@W_6.heap\endcsname{7,19\,s}
\expandafter\def\csname dado@W_6.n\endcsname{10,0\,mi}
\expandafter\def\csname dado@W_6.heapk\endcsname{12,0\,min}
\expandafter\def\csname dado@W_6.vecest\endcsname{16,2\,h}
\expandafter\def\csname dado@W_6.veck\endcsname{67,6\,dias}
\expandafter\def\csname dado@W_6.ratio\endcsname{$\sim$8.100$\times$}
\expandafter\def\csname dado@W_6.heapns\endcsname{3,7}
\expandafter\def\csname dado@W_6.degree\endcsname{14,7}
\expandafter\def\csname dado@W_6.csr\endcsname{1,8\,GB}
\expandafter\def\csname dado@W_6.pico\endcsname{3,3\,GB}
\expandafter\def\csname dado@W_6.matriz\endcsname{900,0\,TB}
\expandafter\def\csname dado@W_6.chainsest\endcsname{7,0\,h}
\expandafter\def\csname dado@W_6.chainsratio\endcsname{$\sim$3.500$\times$}
\expandafter\def\csname dado@W_7.heap\endcsname{31,19\,s}
\expandafter\def\csname dado@W_7.n\endcsname{50,0\,mi}
\expandafter\def\csname dado@W_7.heapk\endcsname{52,0\,min}
\expandafter\def\csname dado@W_7.vecest\endcsname{16,0\,dias}
\expandafter\def\csname dado@W_7.veck\endcsname{4,4\,anos}
\expandafter\def\csname dado@W_7.ratio\endcsname{$\sim$44.000$\times$}
\expandafter\def\csname dado@W_7.heapns\endcsname{11,1}
\expandafter\def\csname dado@W_7.degree\endcsname{2,4}
\expandafter\def\csname dado@W_7.csr\endcsname{1,6\,GB}
\expandafter\def\csname dado@W_7.pico\endcsname{3,4\,GB}
\expandafter\def\csname dado@W_7.matriz\endcsname{22,5\,PB}
\expandafter\def\csname dado@W_7.chainsest\endcsname{7,3\,dias}
\expandafter\def\csname dado@W_7.chainsratio\endcsname{$\sim$20.000$\times$}
\expandafter\def\csname dado@bench.total\endcsname{2,3\,h}
\expandafter\def\csname dado@bench.W7heap\endcsname{52,0\,min}
\expandafter\def\csname dado@bench.vec\endcsname{65,4\,min}
\expandafter\def\csname dado@heapns.growth\endcsname{21$\times$}
\expandafter\def\csname dado@degree.min\endcsname{2,4}
\expandafter\def\csname dado@degree.max\endcsname{122,0}
\expandafter\def\csname dado@matrix.limit\endcsname{66\,mil}
\expandafter\def\csname dado@colab.n\endcsname{722\,mil}
\expandafter\def\csname dado@colab.m\endcsname{2,3\,mi}
\expandafter\def\csname dado@colab.source\endcsname{2722}
\expandafter\def\csname dado@colab.heap\endcsname{112,5\,ms}
\expandafter\def\csname dado@colab.reached\endcsname{65\%}
\expandafter\def\csname dado@cont.origins\endcsname{100}
\expandafter\def\csname dado@W_3.upd\endcsname{4,4}
\expandafter\def\csname dado@W_3.updpct\endcsname{7\%}
\expandafter\def\csname dado@W_3.peakpct\endcsname{95\%}
\expandafter\def\csname dado@W_4.upd\endcsname{3,7}
\expandafter\def\csname dado@W_4.updpct\endcsname{12\%}
\expandafter\def\csname dado@W_4.peakpct\endcsname{92\%}
\expandafter\def\csname dado@W_5.upd\endcsname{1,8}
\expandafter\def\csname dado@W_5.updpct\endcsname{46\%}
\expandafter\def\csname dado@W_5.peakpct\endcsname{61\%}
\expandafter\def\csname dado@W_6.upd\endcsname{2,3}
\expandafter\def\csname dado@W_6.updpct\endcsname{32\%}
\expandafter\def\csname dado@W_6.peakpct\endcsname{75\%}
\expandafter\def\csname dado@W_7.upd\endcsname{1,1}
\expandafter\def\csname dado@W_7.updpct\endcsname{90\%}
\expandafter\def\csname dado@W_7.peakpct\endcsname{14\%}
\expandafter\def\csname dado@colab.giant\endcsname{576.309}
\expandafter\def\csname dado@colab.giant.pct\endcsname{80\%}
\expandafter\def\csname dado@colab.giant.runs\endcsname{81}
\expandafter\def\csname dado@colab.small.runs\endcsname{19}
\expandafter\def\csname dado@colab.small.max\endcsname{8}
\expandafter\def\csname dado@colab.giant.heap\endcsname{138,9\,ms}
\expandafter\def\csname dado@colab.amdahl\endcsname{1,25$\times$}
\expandafter\def\csname dado@colab.kleinberg.d\endcsname{2,71}
\expandafter\def\csname dado@colab.kleinberg.hops\endcsname{9}
\expandafter\def\csname dado@colab.kruskal.d\endcsname{3,48}
\expandafter\def\csname dado@colab.kruskal.hops\endcsname{8}
\expandafter\def\csname dado@colab.tardos.d\endcsname{2,75}
\expandafter\def\csname dado@colab.tardos.hops\endcsname{11}
\expandafter\def\csname dado@colab.figueiredo.d\endcsname{2,94}
\expandafter\def\csname dado@colab.figueiredo.hops\endcsname{8}
\expandafter\def\csname dado@colab.firsttwo.kleinberg\endcsname{52\%}
\expandafter\def\csname dado@W_3.dmin\endcsname{0,80}
\expandafter\def\csname dado@W_3.dmax\endcsname{0,97}
\expandafter\def\csname dado@W_3.hmin\endcsname{9}
\expandafter\def\csname dado@W_3.hmax\endcsname{18}
\expandafter\def\csname dado@W_4.dmin\endcsname{2,02}
\expandafter\def\csname dado@W_4.dmax\endcsname{2,69}
\expandafter\def\csname dado@W_4.hmin\endcsname{13}
\expandafter\def\csname dado@W_4.hmax\endcsname{19}
\expandafter\def\csname dado@W_5.dmin\endcsname{15,43}
\expandafter\def\csname dado@W_5.dmax\endcsname{17,99}
\expandafter\def\csname dado@W_5.hmin\endcsname{12}
\expandafter\def\csname dado@W_5.hmax\endcsname{22}
\expandafter\def\csname dado@W_6.dmin\endcsname{9,60}
\expandafter\def\csname dado@W_6.dmax\endcsname{11,31}
\expandafter\def\csname dado@W_6.hmin\endcsname{11}
\expandafter\def\csname dado@W_6.hmax\endcsname{22}
\expandafter\def\csname dado@W_7.dmin\endcsname{82,25}
\expandafter\def\csname dado@W_7.dmax\endcsname{100,69}
\expandafter\def\csname dado@W_7.hmin\endcsname{32}
\expandafter\def\csname dado@W_7.hmax\endcsname{46}


\pgfplotsset{every axis/.append style={tick label style={font=\footnotesize},
  label style={font=\small}, legend style={font=\footnotesize}}}
\tikzset{
  vert/.style={circle, draw, thick, minimum size=16pt, inner sep=0pt, font=\small},
  explorado/.style={vert, fill=gray!30},
  peso/.style={font=\scriptsize, fill=white, inner sep=1pt},
  cel/.style={draw, minimum width=0.66cm, minimum height=0.5cm, inner sep=0pt,
    font=\scriptsize, anchor=west},
  rot/.style={font=\scriptsize, anchor=east, inner sep=2pt}}

\begin{document}

% ── capa (sem número; a numeração começa na página seguinte) ──
\begin{titlepage}
\linespread{1.0}\selectfont
\centering
\includegraphics[width=3.2cm]{../comum/minerva.png}\par\vspace{1.2em}
{\large Universidade Federal do Rio de Janeiro}\par\vspace{0.3em}
Escola Politécnica\par\vspace{0.3em}
COS242 Teoria dos Grafos, 2026/2\par
\vfill
{\Large Trabalho de Disciplina, Parte 2}\par\vspace{1.2em}
{\huge\bfseries Caminhos mínimos em grafos com pesos:\par\vspace{0.3em}
 Dijkstra com vetor e com heap\par}
\vfill
{\large Felipe Marroni\par\vspace{0.3em} Hygor Knust\par}
\vspace{2.5em}
Professor: Daniel Ratton Figueiredo\par
\vfill
{\small\url{https://github.com/knusth/cos-242-grafos}}\par\vspace{0.8em}
Rio de Janeiro, \today
\end{titlepage}
\setcounter{page}{1}

\begin{quote}\small
\textbf{Resumo.} Estendemos a biblioteca de grafos para grafos com pesos e implementamos o
algoritmo de Dijkstra uma única vez, de modo que ele pode guardar as distâncias em duas estruturas,
um vetor ou um heap. Nos grafos \W{3} a \W{7} do site
da disciplina, o heap foi de \dado{W_3.ratio} a \dado{W_7.ratio} mais rápido que o vetor: no
\W{7}, uma execução levou \dado{W_7.heap} com o heap e levaria \dado{W_7.vecest} com o vetor.
O número de atualizações no heap foi de \dado{W_3.updpct} a \dado{W_7.updpct} de $m$: o termo
$m\log n$ da análise é pessimista nos grafos mais densos.
\end{quote}

\section{Introdução}

A biblioteca agora aceita grafos com pesos e calcula a distância e o caminho mínimo de um
vértice a todos os outros com o algoritmo de Dijkstra, numa única implementação que pode usar
um vetor ou um heap para guardar as distâncias.

\section{Decisões de projeto e implementação}
\label{sec:projeto}

\subsection{Representação do grafo com pesos}

A biblioteca guarda o grafo em vetores de adjacência (aula~3) concatenados num vetor de
\emph{vizinhos}, com um vetor de \emph{início} que indica onde começa a lista de cada vértice
(Figura~\ref{fig:repr}); cada aresta aparece nas listas das duas extremidades. Para os pesos,
acrescentamos um terceiro vetor, alinhado ao de vizinhos: ao percorrer os vizinhos de $u$,
Dijkstra lê o peso de cada aresta na mesma posição, sem nenhuma busca.

\begin{figure}[!htb]
\centering
\begin{tikzpicture}
% grafo de exemplo do enunciado
\node[vert] (1) at (0,1.5) {1};
\node[vert] (2) at (1.3,2.5) {2};
\node[vert] (5) at (1.6,0.9) {5};
\node[vert] (3) at (3.1,1.9) {3};
\node[vert] (4) at (2.9,0.0) {4};
\draw (1) -- node[peso] {0,1} (2);
\draw (2) -- node[peso] {0,2} (5);
\draw (5) -- node[peso] {5} (3);
\draw[ruim, very thick] (3) -- node[peso, text=ruim] {$-9{,}5$} (4);
\draw (4) -- node[peso] {2,3} (5);
\draw (1) -- node[peso] {1} (5);
% vetores
\def\xa{5.6}
\node[rot] at (\xa,2.4) {início};
\node[rot] at (\xa,1.0) {vizinhos};
\node[rot] at (\xa,0.5) {pesos};
\foreach \v [count=\i from 0] in {0,2,4,6,8,12} {
  \node[cel] (o\i) at (\xa+0.66*\i,2.4) {\v};
}
\foreach \i in {1,...,5} {
  \node[font=\scriptsize, gray] at ($(o\the\numexpr\i-1\relax.north)+(0,0.18)$) {\i};
}
\node[font=\scriptsize, gray] at ($(o5.north)+(0,0.18)$) {fim};
\foreach \v [count=\i from 0] in {2,5,1,5,5,4,3,5,2,3,4,1} {
  \node[cel] (n\i) at (\xa+0.66*\i,1.0) {\v};
  \node[font=\tiny, gray] at ($(n\i.south)+(0,-0.68)$) {\i};
}
\foreach \p [count=\i from 0] in {{0,1},1,{0,1},{0,2},5,{$-9{,}5$},{$-9{,}5$},{2,3},{0,2},5,{2,3},1} {
  \node[cel] (p\i) at (\xa+0.66*\i,0.5) {\p};
}
% blocos de cada vértice
\foreach \a/\b/\u in {0/1/1,2/3/2,4/5/3,6/7/4,8/11/5} {
  \draw[decorate, decoration={brace, amplitude=3pt}, gray]
    ($(n\a.north west)+(0.04,0.05)$) -- node[above=1pt, font=\tiny, gray] {de \u}
    ($(n\b.north east)+(-0.04,0.05)$);
}
\draw[-{Stealth[length=4pt]}, heapc, thick] (o4.south) to[out=-90, in=90] ($(n8.north)+(0,0.42)$);
\end{tikzpicture}
\caption{Grafo de exemplo do enunciado e sua representação. Os vizinhos de $u$ ocupam as
posições de $\ini[u]$ a $\ini[u+1]-1$ (numeradas a partir de 0, abaixo); a seta mostra que os
vizinhos de 5 começam na posição 8. Por causa da
aresta $(3,4)$, de peso negativo, a biblioteca não executa Dijkstra neste grafo.}
\label{fig:repr}
\end{figure}

Essa representação ocupa 4 bytes por vértice e 24 bytes por aresta (4 do vizinho e 8 do peso,
em cada sentido): \dado{W_3.csr} no \W{3} e \dado{W_7.csr} no \W{7}. A matriz de adjacência,
que a biblioteca oferece como opção de compilação, guarda para cada par de vértices um byte que
indica a aresta e um peso de 8 bytes, ou seja, $9n^2$ bytes: \dado{W_3.matriz} no \W{3} e
\dado{W_7.matriz} no \W{7}. Com os \dado{ram} de memória da máquina usada, a matriz só caberia
até cerca de \dado{matrix.limit} vértices, menos que o menor grafo do trabalho.

\textbf{Pesos negativos.} Ao ler o arquivo, a biblioteca registra se algum peso é negativo.
Nesse caso, ela não executa Dijkstra e informa que ``ainda não implementa caminhos
mínimos com pesos negativos''. É o que acontece com o grafo da Figura~\ref{fig:repr}. Arestas repetidas
no arquivo são mantidas; como Dijkstra só atualiza $\dist[v]$ quando encontra um valor menor,
prevalece a de menor peso.

\subsection{Dijkstra com vetor e com heap}

A cada passo, o algoritmo da aula~8 extrai, dentre os vértices não explorados ($V-S$), o de
menor $\dist[u]$, o adiciona ao conjunto $S$ dos explorados e, para cada vizinho $v$, atualiza
$\dist[v]$ se $\dist[v] > \dist[u] + w(u,v)$. As duas versões só diferem nessas duas operações,
\emph{extrair} e \emph{atualizar}; por isso, escrevemos o laço uma única vez, como um
\emph{template} parametrizado pela fila de prioridades. A Figura~\ref{fig:passo} mostra as duas
estruturas no mesmo passo do algoritmo.

\textbf{Vetor.} A estrutura guarda a distância de $v$ na posição $v$. Atualizar é escrever
nessa posição, em $O(1)$; extrair o mínimo exige percorrer as $n$ posições, em $O(n)$. Como cada
vértice alcançado é extraído uma vez, o total num grafo conexo é $O(n^2)$, mais $O(m)$ para as
atualizações.

\textbf{Heap.} A estrutura é um heap binário de mínimo, como na aula~9, em que a chave de cada
vértice é sua distância. Extrair o mínimo custa $O(\log n)$.
Atualizar a distância de um vértice que já está no heap exige saber onde ele está; para isso,
mantemos um vetor de posições, em que $\pos[v]$ é a posição de $v$ no heap, corrigida a cada
troca de elementos. Com ele, a atualização vai direto ao elemento de $v$, reduz sua chave e o
faz subir até a posição correta, também em $O(\log n)$. São no máximo $n$ extrações e $m$
atualizações, e o total é $O((n+m)\log n)$. Sem o vetor de posições, a alternativa seria inserir
uma nova cópia do vértice a cada atualização e descartar as cópias antigas ao extraí-las; o heap
poderia então chegar a $m$ elementos.

A Listagem~\ref{lst:estruturas} mostra como essas estruturas estão representadas no código,
com $n$ vértices e $m$ arestas.

\begin{lstlisting}[float=!htb, caption={Estruturas de dados do grafo e das duas filas de
prioridades (campos simplificados da biblioteca).}, label=lst:estruturas]
// Grafo: vizinhos de u em neighbors_[offsets_[u] .. offsets_[u+1]-1]
uint32 offsets_[n+1];      // início de cada lista              4(n+1) bytes
uint32 neighbors_[2m];     // cada aresta nas duas listas           8m bytes
double weights_[2m];       // peso do vizinho na mesma posição     16m bytes

// Dijkstra
double distance[n];        // dist[]: (*$\infty$*) até ser alcançado
uint32 parent[n];          // pai[]: árvore geradora induzida pela busca

// Fila com vetor
double estimate_[n];       // (*$\infty$*) = fora da fila; pop() percorre as n posições

// Fila com heap
struct Entry { double key; uint32 vertex; };
Entry  heap_[n];           // filhos de i em 2i+1 e 2i+2 (na aula 9, 2i e 2i+1)
uint32 position_[n];       // pos[v]: índice de v em heap_, ou nenhum
uint32 size_;
\end{lstlisting}

\begin{figure}[!htb]
\centering
\begin{minipage}[b]{0.34\linewidth}\centering
\begin{tikzpicture}[scale=0.92]
\node[explorado] (1) at (0,0) {1};
\node[explorado] (2) at (1.3,0.9) {2};
\node[vert] (3) at (1.3,-0.9) {3};
\node[vert] (4) at (2.8,0.9) {4};
\node[vert] (5) at (2.8,-0.9) {5};
\node[vert] (6) at (4.1,0) {6};
\begin{scope}[gray!60]
\draw (1) -- node[peso] {5} (3);
\draw (3) -- node[peso] {3} (5);
\draw (4) -- node[peso] {1} (5);
\draw (4) -- node[peso] {1} (6);
\draw (5) -- node[peso] {2} (6);
\end{scope}
\draw[very thick] (1) -- node[peso] {2} (2);
\draw[very thick] (2) -- node[peso] {1} (3);
\draw[very thick] (2) -- node[peso] {4} (4);
\draw[very thick] (2) -- node[peso, pos=0.35] {7} (5);
\begin{scope}[font=\scriptsize\bfseries, heapc]
\node[above=1pt of 1] {0};
\node[above=1pt of 2] {2};
\node[below=1pt of 3] {3};
\node[above=1pt of 4] {6};
\node[below=1pt of 5] {9};
\node[above=1pt of 6] {$\infty$};
\end{scope}
\end{tikzpicture}\\[2pt]
{\small (a) grafo, $S=\{1,2\}$}
\end{minipage}%
\begin{minipage}[b]{0.3\linewidth}\centering
\begin{tikzpicture}
\foreach \v/\d [count=\i from 0] in {1/0,2/2,3/3,4/6,5/9,6/$\infty$} {
  \ifnum\i<2
    \node[cel, minimum width=0.55cm, fill=gray!30] (c\i) at (0.55*\i,0) {\d};
  \else
    \node[cel, minimum width=0.55cm] (c\i) at (0.55*\i,0) {\d};
  \fi
  \node[font=\scriptsize, gray] at ($(c\i.north)+(0,0.18)$) {\v};
}
\draw[heapc, very thick] (c2.north west) rectangle (c2.south east);
\draw[-{Stealth[length=4pt]}, gray] ($(c0.south west)+(0,-0.25)$) -- node[below, font=\scriptsize,
  text=black] {percorre as $n$ posições} ($(c5.south east)+(0,-0.25)$);
\node at (0,1.0) {};
\end{tikzpicture}\\[2pt]
{\small (b) vetor de distâncias}
\end{minipage}%
\begin{minipage}[b]{0.34\linewidth}\centering
\begin{tikzpicture}
\node[draw, rounded corners=2pt, font=\scriptsize, fill=heapc!12] (h1) at (0,1.5) {3: 3};
\node[draw, rounded corners=2pt, font=\scriptsize] (h2) at (-0.8,0.8) {4: 6};
\node[draw, rounded corners=2pt, font=\scriptsize] (h3) at (0.8,0.8) {5: 9};
\draw (h1) -- (h2); \draw (h1) -- (h3);
\node[font=\tiny, gray, left=1pt of h1] {1};
\node[font=\tiny, gray, left=1pt of h2] {2};
\node[font=\tiny, gray, right=1pt of h3] {3};
\node[rot] at (-1.0,-0.1) {$\pos$};
\foreach \p [count=\i from 1] in {---,---,1,2,3,---} {
  \node[cel, minimum width=0.45cm] (q\i) at (-1.0+0.45*\i-0.45,-0.1) {\p};
  \node[font=\tiny, gray] at ($(q\i.south)+(0,-0.15)$) {\i};
}
\end{tikzpicture}\\[2pt]
{\small (c) heap e vetor de posições}
\end{minipage}
\caption{Estado de Dijkstra a partir do vértice 1, depois de explorar 1 e 2 (em cinza). (a) Em
verde, $\dist$ de cada vértice; as arestas em preto ligam cada vértice ao seu pai. (b) Com o
vetor, extrair o mínimo percorre todas as posições, ignorando as de $S$, e encontra o vértice 3.
(c) Com o heap, o mínimo está na raiz (cada elemento é ``vértice: chave''; posições a partir de
1, como na aula~9), e $\pos$ indica onde está cada vértice (--- se não está no heap). No passo
seguinte, 3 é explorado e $\dist[5]$ cai de 9 para $3+3=6$: no vetor, basta escrever na posição
5; no heap, $\pos[5]$ leva ao elemento do vértice 5, que sobe se for preciso.}
\label{fig:passo}
\end{figure}

Dois detalhes diferem do pseudocódigo da aula. O laço termina quando a fila de prioridades fica
vazia, e não quando $S=V$: num grafo desconexo, os vértices fora da componente da origem nunca
entram na fila e ficam com $\dist=\infty$. E, sempre que $\dist[v]$ é atualizada a partir de
$u$, guardamos $\pai[v]=u$ (aula~5), de onde reconstruímos os caminhos.

\section{Metodologia}
\label{sec:metodo}

Rodamos todos os experimentos num \dado{cpu} com \dado{ram} de memória, compilando com
\dado{compiler} e \cod{-O3}. Em cada grafo, sorteamos $k=\dado{k}$ vértices de origem (semente
\dado{seed}) e usamos os mesmos nas duas versões. Medimos apenas a execução de Dijkstra, sem a
leitura do arquivo. Para cada origem executada
com as duas versões, conferimos que elas alcançaram os mesmos vértices com a mesma soma de
distâncias. Para contar as operações do heap, repetimos as mesmas origens, fora das medições de
tempo, com um contador de atualizações e do tamanho do heap. Outros processos rodavam na máquina durante parte das medições, o que explica parte
do desvio-padrão da Tabela~\ref{tab:tempos}.

Com o vetor, as \dado{k} origens levariam \dado{W_4.veck} no \W{4} e \dado{W_7.veck} no \W{7}.
Por isso, no \W{4} executamos o vetor a partir de apenas \dado{val.W_4.runs} das origens, e do
\W{5} ao \W{7} projetamos o tempo. Como cada extração percorre as mesmas $n$ posições, todas
custam aproximadamente o mesmo. Executamos o vetor a partir de 5 origens, interrompemos cada
execução depois de algumas extrações (\dado{pops.limit.W_5} no \W{5}, \dado{pops.limit.W_7} no \W{7}),
medimos o tempo médio de uma extração e o multiplicamos pelo número de vértices alcançados, que
já conhecíamos das execuções com heap. No \W{3} e no \W{4}, onde executamos o vetor até o fim, a projeção
diferiu da mediana medida em \dado{val.W_3.err} e \dado{val.W_4.err}, respectivamente.

\section{Estudos de caso}
\label{sec:casos}

\subsection{Distância e caminho mínimo a partir do vértice 10}

A Tabela~\ref{tab:dist} traz a distância do vértice 10 aos vértices 20, 30, 40, 50 e 60 em cada
grafo, com o número de arestas do caminho mínimo; os caminhos completos estão no
Apêndice~\ref{ap:caminhos}. As distâncias crescem à medida que o grau médio diminui: ficam entre
\dado{W_3.dmin} e \dado{W_3.dmax} no \W{3}, de grau médio \dado{W_3.degree}, e entre
\dado{W_7.dmin} e \dado{W_7.dmax} no \W{7}, de grau médio \dado{W_7.degree}. O caminho mínimo também
usa mais arestas: de
\dado{W_7.hmin} a \dado{W_7.hmax} no \W{7}. A comparação entre \W{5} e \W{6}, que têm o mesmo
número de vértices, mostra o mesmo efeito: no \W{6}, com quase o dobro de arestas, todas as
distâncias são menores.

\begin{table}[!htb]
\centering\small
\caption{Distância $d$ do vértice 10 ao vértice $t$ e número de arestas do caminho mínimo.
Todos os destinos são alcançáveis.}
\label{tab:dist}
\setlength{\tabcolsep}{5pt}
\begin{tabular}{@{}r*{5}{rr}@{}}
\toprule
 & \multicolumn{2}{c}{\W{3}} & \multicolumn{2}{c}{\W{4}} & \multicolumn{2}{c}{\W{5}} &
   \multicolumn{2}{c}{\W{6}} & \multicolumn{2}{c}{\W{7}} \\
\cmidrule(lr){2-3}\cmidrule(lr){4-5}\cmidrule(lr){6-7}\cmidrule(lr){8-9}\cmidrule(l){10-11}
$t$ & $d$ & arestas & $d$ & arestas & $d$ & arestas & $d$ & arestas & $d$ & arestas \\
\midrule
\tabelaDistancias
\bottomrule
\end{tabular}
\end{table}

\subsection{Tempo médio das duas versões}

\begin{table}[!htb]
\centering\small
\caption{Tempo médio de uma execução de Dijkstra, de uma origem a todos os vértices; grau médio
$=2m/n$. Heap: média $\pm$ desvio-padrão nas \dado{k} origens. Vetor: média de
\dado{val.W_3.runs} execuções no \W{3} e de \dado{val.W_4.runs} no \W{4}; nos demais grafos, em
itálico, projeção (Seção~\ref{sec:metodo}). As razões marcadas com $\sim$ vêm da projeção e têm
dois algarismos significativos. Última coluna: tempo do heap, em ns, dividido por
$(n+m)\log_2 n$; seria constante se toda operação custasse o mesmo.}
\label{tab:tempos}
\setlength{\tabcolsep}{6pt}
\begin{tabular}{@{}lrrrrrrr@{}}
\toprule
grafo & $n$ & $m$ & grau médio & heap & vetor & vetor/heap & ns/op.\ heap \\
\midrule
\tabelaTempos
\bottomrule
\end{tabular}
\end{table}

A Tabela~\ref{tab:tempos} e a Figura~\ref{fig:tempos} comparam o tempo médio das duas versões.
O heap foi mais rápido em todos os grafos, por um fator que foi de \dado{W_3.ratio} no \W{3} a
\dado{W_7.ratio} no \W{7}, como prevê a análise de complexidade: nesses grafos esparsos, o grau
médio não passa de \dado{degree.max}, e $(n+m)\log n$ fica muito abaixo de $n^2$. A razão
cresceu com $n$ porque o tempo do vetor cresce com $n^2$, enquanto o do heap cresce pouco mais
que linearmente. Com o vetor, as \dado{k} origens do estudo de caso levariam \dado{W_7.veck} no
\W{7}; com o heap, levaram \dado{W_7.heapk}. Só a versão com heap tornou o estudo viável em todos
os grafos.

\begin{figure}[!htb]
\centering
\begin{tikzpicture}
\begin{axis}[grafosW, ybar, bar width=9pt, ymode=log, log origin=infty,
  width=0.78\linewidth, height=6.2cm, ymin=0.01, ymax=4e7,
  ytick={0.01,1,60,3600,86400,2592000},
  yticklabels={10\,ms,1\,s,1\,min,1\,h,1\,dia,30\,dias},
  ymajorgrids, grid style={gray!18}, enlarge x limits=0.14, xtick=data,
  ylabel={tempo por origem},
  legend style={at={(0.02,0.98)}, anchor=north west, draw=none, fill=white,
    legend cell align=left}]
\addplot[fill=heapc, draw=none, bar shift=-5pt] coordinates {\barrasHeap};
\addplot[fill=ruim, draw=none, bar shift=5pt] coordinates {\barrasVetor};
\addplot[fill=ruim!18, draw=ruim, densely dashed, bar shift=5pt] coordinates {\barrasVetorProj};
\legend{heap, vetor (medido), vetor (projeção)}
\razoesW
\end{axis}
\end{tikzpicture}
\caption{Tempo de uma execução de Dijkstra por origem, em escala logarítmica. As barras
tracejadas são projeções; o número sobre cada par é a razão vetor/heap.}
\label{fig:tempos}
\end{figure}

\textbf{Quantas operações o heap faz.} A análise conta até $m$ atualizações, uma por aresta, cada
uma em $O(\log n)$. Contamos as operações nas \dado{cont.origins} origens
(Figura~\ref{fig:operacoes}). As atualizações de $\dist$ ficaram muito abaixo de $m$: foram
\dado{W_3.upd} por vértice no \W{3}, onde o limite é $m/n=61$, e \dado{W_7.upd} no \W{7}; em
proporção a $m$, foram de \dado{W_3.updpct} no \W{3} a \dado{W_7.updpct} no \W{7}. O número de
atualizações por vértice quase não cresceu com o grau, porque a maioria dos testes
$\dist[v] > \dist[u] + w(u,v)$ falha: $v$ já foi alcançado por um caminho mais curto. Já o termo
$\log n$ não é folgado: nos grafos densos, o heap passou a conter quase todos os vértices não
explorados logo no início (pico de \dado{W_3.peakpct} de $n$ no \W{3}); no \W{7}, o mais
esparso, o pico foi \dado{W_7.peakpct} de $n$, o que reduziu a altura do heap em apenas três
níveis. Assim, o custo do heap nestes grafos fica mais perto de $O(m + n\log n)$ que de
$O((n+m)\log n)$: as $m$ arestas são examinadas, mas, nos grafos mais densos, só uma pequena
fração delas gera operação no heap. Isso também explica o \W{5} e o \W{6}, que têm o mesmo $n$: o vetor levaria o mesmo tempo
nos dois (\dado{W_5.vecest} e \dado{W_6.vecest}), mas o heap fez \dado{W_6.upd} atualizações
por vértice no \W{6}, contra \dado{W_5.upd} no \W{5}, e a razão caiu de \dado{W_5.ratio} para
\dado{W_6.ratio}.

\begin{figure}[!htb]
\centering
\begin{minipage}[b]{0.49\linewidth}\centering
\begin{tikzpicture}
\begin{axis}[width=\linewidth, height=5.4cm, xmode=log, ymode=log,
  xlabel={grau médio}, ylabel={atualizações por vértice},
  xmin=1.5, xmax=250, ymin=0.55, ymax=130, ymajorgrids, grid style={gray!18},
  xtick={2,10,100}, xticklabels={2,10,100}, ytick={1,10,100}, yticklabels={1,10,100}]
\addplot[gray, dashed, domain=1.5:250, samples=2] {x/2};
\node[font=\scriptsize, gray, rotate=33, anchor=south] at (axis cs:30,15) {limite $m/n$};
\addplot[only marks, mark=*, heapc] coordinates {\pontosAtualizacoes};
\rotulosAtualizacoes
\end{axis}
\end{tikzpicture}\\[2pt]
{\small (a) atualizações de $\dist$}
\end{minipage}\hfill
\begin{minipage}[b]{0.49\linewidth}\centering
\begin{tikzpicture}
\begin{axis}[width=\linewidth, height=5.4cm, xmin=0, xmax=1, ymin=0, ymax=1,
  xlabel={vértices extraídos / $n$}, ylabel={tamanho do heap / $n$},
  ymajorgrids, grid style={gray!18}, legend style={draw=none, fill=white, font=\scriptsize,
  at={(0.98,0.98)}, anchor=north east, legend cell align=left}]
\addplot[gray, dashed, domain=0:1, samples=2, forget plot] {1-x};
\addplot[black, thick] coordinates {\csname heapTam@W_3\endcsname}; \addlegendentry{\W{3}}
\addplot[gray, thick] coordinates {\csname heapTam@W_4\endcsname}; \addlegendentry{\W{4}}
\addplot[heapc, thick] coordinates {\csname heapTam@W_5\endcsname}; \addlegendentry{\W{5}}
\addplot[heapc!50, thick] coordinates {\csname heapTam@W_6\endcsname}; \addlegendentry{\W{6}}
\addplot[ruim, thick] coordinates {\csname heapTam@W_7\endcsname}; \addlegendentry{\W{7}}
\end{axis}
\end{tikzpicture}\\[2pt]
{\small (b) tamanho do heap durante uma execução}
\end{minipage}
\caption{Operações do heap. (a) Média de atualizações de $\dist$ por vértice nas
\dado{cont.origins} origens, contra o grau médio; a reta tracejada é o limite da análise, $m$
atualizações. (b) Tamanho do heap ao longo de uma execução a partir de uma das origens; a reta
tracejada, $n$ menos os extraídos, é o máximo possível.}
\label{fig:operacoes}
\end{figure}

\textbf{Por que o heap fica mais lento por operação.} O tempo do heap cresceu mais rápido que
$(n+m)\log n$: dividido por essa quantidade (última coluna da Tabela~\ref{tab:tempos}), ele subiu
de \dado{W_3.heapns}\,ns no \W{3} para \dado{W_7.heapns}\,ns no \W{7}, \dado{heapns.growth}
mais. A análise de complexidade supõe que toda leitura custa o mesmo, mas o custo depende do nível
da hierarquia de memória em que o dado está:

\begin{center}\small\setlength{\tabcolsep}{6pt}
\begin{tabular}{@{}lrrl@{}}
\toprule
memória & tamanho & tempo de leitura & o vetor de distâncias ($8n$ bytes) cabe \\
\midrule
cache L1 (por núcleo) & 128\,KiB & \dado{lat.l1}\,ns & se $n\le16$ mil: nenhum grafo \\
cache L2 (compartilhada) & 16\,MiB & \dado{lat.l2}\,ns & se $n\le2$ milhões: \W{3} e \W{4} \\
RAM & \dado{ram} & \dado{lat.dram}\,ns & sempre; \W{5} a \W{7} só aqui \\
\bottomrule
\end{tabular}
\end{center}

\noindent A cada aresta examinada, Dijkstra lê $\dist[v]$ de um vizinho $v$, que pode estar em
qualquer posição do vetor de distâncias. A Figura~\ref{fig:latencia} mostra o custo dessa leitura conforme o tamanho
dos dados. O vetor de distâncias do \W{3} e do \W{4} cabe na L2, e o do \W{5} em diante não, o que é
compatível com o salto da última coluna da Tabela~\ref{tab:tempos} entre o \W{4}
(\dado{W_4.heapns}\,ns) e o \W{5} (\dado{W_5.heapns}\,ns). A versão com vetor não sofre esse
efeito, pois percorre o vetor de distâncias sequencialmente.

\begin{figure}[!htb]
\centering
\begin{tikzpicture}
\begin{axis}[width=0.72\linewidth, height=5.4cm, xmode=log, ymode=log,
  xlabel={tamanho dos dados percorridos}, ylabel={ns por leitura},
  ymajorgrids, grid style={gray!18}, xmin=1e4, xmax=4e9, ymin=0.7, ymax=200,
  xtick={16384,1048576,67108864,1073741824}, xticklabels={16\,KiB,1\,MiB,64\,MiB,1\,GiB},
  ytick={1,10,100}, yticklabels={1,10,100}]
\draw[gray, dashed] (axis cs:131072,0.7) -- (axis cs:131072,200);
\draw[gray, dashed] (axis cs:16777216,0.7) -- (axis cs:16777216,200);
\node[font=\scriptsize, gray, anchor=north east] at (axis cs:131072,180) {L1};
\node[font=\scriptsize, gray, anchor=north east] at (axis cs:16777216,180) {L2};
\addplot[black, thick, mark=*, mark size=1.2pt] coordinates {\latencia};
\marcasW
\end{axis}
\end{tikzpicture}
\caption{Tempo de uma leitura numa posição imprevisível da memória, conforme o tamanho dos
dados percorridos; eixos em escala logarítmica. As linhas tracejadas marcam o tamanho das caches L1 e
L2, e as marcas no eixo horizontal, o tamanho do vetor de distâncias de cada grafo.}
\label{fig:latencia}
\end{figure}

\subsection{Rede de colaboração}

A Figura~\ref{fig:arvore} mostra a distância e o caminho mínimo de Edsger W.~Dijkstra até Jon M.~Kleinberg,
J.\,B.~Kruskal, Éva Tardos e Daniel R.~Figueiredo. Alan M.~Turing não aparece: está em outra componente conexa, e a distância
até ele é infinita.

\begin{figure}[!htb]
\centering
\begin{tikzpicture}[
  pessoa/.style={font=\footnotesize,fill=white,inner sep=1.5pt},
  raiz/.style={font=\small\bfseries,fill=white,inner sep=2pt},
  folha/.style={font=\footnotesize,fill=gray!10,draw=gray!60!black,rounded corners,align=center,inner sep=2pt},
  peso/.style={font=\scriptsize,inner sep=1pt,anchor=west,text=gray!70!black,fill=white},
  barato/.style={gray!60},
  caro/.style={ruim,very thick}]
\arvoreColab
\end{tikzpicture}
\caption{Caminhos mínimos a partir de Edsger W.~Dijkstra. Cada nível desce uma aresta; as
folhas trazem a distância. Os pesos estão arredondados para duas casas; por isso, a soma pode
diferir da distância na última casa. Em vermelho, as arestas de peso 1,00, as mais altas destes
caminhos.}
\label{fig:arvore}
\end{figure}

O caminho mínimo nem sempre é o de menos arestas: Kleinberg está a \dado{colab.kleinberg.d}, por um
caminho de \dado{colab.kleinberg.hops} arestas, e Kruskal, a \dado{colab.kruskal.d}, por
\dado{colab.kruskal.hops} arestas. O caminho até Kruskal passa por duas arestas de peso 1,00, em
vermelho na figura, e o até Kleinberg, por uma.

\section{Conclusão}

Nos estudos de caso, o heap foi mais
rápido que o vetor em todos os grafos, de \dado{W_3.ratio} a \dado{W_7.ratio}, e foi a única
versão viável a partir do \W{4}. A diferença cresce com $n$ e diminui quando $m$ cresce com $n$
fixo, como prevê a análise de complexidade. As contagens mostram que, nos grafos mais densos, só uma fração
das arestas gera atualização no heap, de modo que o termo $m\log n$ da análise é pessimista. Na
rede de colaboração, o caminho mínimo nem sempre é o de menos arestas. Como trabalho futuro,
falta calcular caminhos mínimos com pesos negativos, que a biblioteca hoje recusa.


\appendix
\section{Caminhos mínimos a partir do vértice 10}
\label{ap:caminhos}

{\small\linespread{1.0}\selectfont\setlength{\tabcolsep}{5pt}
\begin{longtable}{@{}l r r r >{\ttfamily\scriptsize\raggedright\arraybackslash}p{0.66\linewidth}@{}}
\caption{Caminhos mínimos do vértice 10 a cada destino $t$ da Tabela~\ref{tab:dist}, com os
vértices numerados como nos arquivos.}\label{tab:caminhos}\\
\toprule
grafo & $t$ & $d$ & arestas & \multicolumn{1}{l}{caminho} \\
\midrule
\endfirsthead
\toprule
grafo & $t$ & $d$ & arestas & \multicolumn{1}{l}{caminho} \\
\midrule
\endhead
\bottomrule
\endlastfoot
\tabelaCaminhos
\end{longtable}}

\end{document}
